\section{Discussion and conclusion}
\label{sec:conclusion}

All five layouts follow from one principle: fictitious sources must enclose the singularities of the continued scattered field, or reach them where they meet $\Gamma$. Any non-resonant inner radius works on the sphere, whereas the focal ring bounds the torus surface from below, a focal curve runs within $6\%$ of the almond rim, and on the cube and tetrahedron the singular set meets $\Gamma$ along the edges and vertices of the skeleton. Wherever it touches $\Gamma$, singular tails with analytical wedge and conical exponents restore rapid convergence with few unknowns. We note two main limitations.

\emph{Generality and accuracy.} Each layout was built from the singular set of its geometry rather than generated automatically; \autoref{thm:skeleton} applies to convex conducting polyhedra only, while the torus bound \eqref{eq:torusradius} and almond focal curve rest on analytical heuristics and numerical tests rather than three-dimensional Maxwell proofs. Graded Bempp confirms the field to $6\times10^{-5}$ on the cube ($2\times10^{-4}$ on the regular tetrahedron and $1.2\times10^{-3}$ without symmetry on the asymmetric one, where fewer unknowns fit in $8$\,GB); beyond that, self-convergence places our cube reference error near $5\times10^{-6}$, with a fixed-$q$ floor further down.

\emph{Electrical size and comparison.} All bodies here are at most nine wavelengths long, dense factorisation scales cubically, and our timings compare against a lowest-order dense boundary element solver rather than a high-order or fast-multipole code. Because the skeleton columns separate into a coarse interior core and distinct edge walls with moderate condition numbers (below $3\times10^7$, \autoref{sec:method}), however, iterative least squares works without forming $A$. Orthogonalising each edge wall against the coarse core (inner points, vertex rays and two coarsest levels) stores $11\text{--}14$ times less memory than $A$ and converges to relative normal-equation residual $10^{-4}$ in $17\text{--}22$ matrix-free panel-tree LSMR \cite{fong2011} iterations on the cube ($J = 4,6,8$ at $k = 2\pi$ and $J = 6$ at $k = 4\pi$; $7\text{--}53$\,MiB versus $88\text{--}625$\,MiB) and $33\text{--}35$ on the asymmetric tetrahedron ($J = 4,5$; $158\text{--}259$\,MiB versus $2.10\text{--}3.51$\,GiB), matching direct solves to $1\text{--}6\times10^{-4}$ (whereas $54\text{--}100$ diagonal-scaled iterations leave $3\times10^{-3}\text{--}6\times10^{-2}$); fast multipoles would speed each product. Adding four and twelve modes per unit length handles $k = 4\pi$ and $6\pi$ (\autoref{sec:cubeother}), and a preliminary $k = 8\pi$ test reaches $10^{-3}$ with $24$ extra modes and five nodes per panel (six for finer levels).

In conclusion, on convex polyhedra and the NASA almond, the Maxwell MFS in three dimensions converges once sources sit on an interior set reaching the surface singularities, cluster towards its boundary contact, and continue below cut-off via parameter-free singular tails. With these layouts, root-exponential convergence holds down to $2\times10^{-5}$ on the cube and $10^{-4}$ on the regular tetrahedron (and without symmetry on an asymmetric one), and the NASA almond radar cross-section matches a published reference within $0.12$\,dB on average up to $10.25$\,GHz. Automating layout construction for general geometries is the natural next step.


